% Generated draft: requires mathematical and source review.
\documentclass[11pt]{article}
\usepackage[margin=25mm]{geometry}
\usepackage{fontspec}
\setmainfont{Latin Modern Roman}
\usepackage{amsmath,amssymb,mathtools,unicode-math}
\setmathfont{Latin Modern Math}
\usepackage{enumitem,needspace,xurl,hyperref}
\setlength{\parindent}{0pt}
\setlength{\parskip}{5pt}
\setlength{\emergencystretch}{3em}
\newcommand{\N}{\mathbb{N}}
\newcommand{\Z}{\mathbb{Z}}
\newcommand{\Q}{\mathbb{Q}}
\newcommand{\R}{\mathbb{R}}
\newcommand{\C}{\mathbb{C}}
\newcommand{\F}{\mathbb{F}}
\newcommand{\A}{\mathbb{A}}
\newcommand{\PP}{\mathbb{P}}
\newcommand{\E}{\mathbb{E}}
\newcommand{\eps}{\varepsilon}
\newcommand{\dd}{\,\mathrm{d}}
\begin{document}
\section*{UnsolvedMath ID: 700609}
\textbf{Source problem code:} catalogue entry
\subsection{Definitions and mathematical statement}
Let \(p\) be a prime, let \(k\) be any field of characteristic \(p\), and let \(r\geq 2\) be an integer. Let
\[
F=k\langle x_1,\ldots,x_r\rangle_{+}
\]
be the free associative \(k\)-algebra without identity consisting of polynomials with zero constant term in the indicated noncommuting variables. Let \(I\) be a two-sided ideal of \(F\), and put \(A=F/I\). Assume that \(A\) is nil but not nilpotent: every \(a\in A\) satisfies \(a^m=0\) for some integer \(m\geq 1\) depending on \(a\), while for every integer \(n\geq 1\) there exist \(a_1,\ldots,a_n\in A\) with \(a_1\cdots a_n\neq 0\).

In the algebra obtained by adjoining an identity to \(A\), the adjoint group \(1+A\) consists of the elements \(1+a\), with multiplication
\[
(1+a)(1+b)=1+(a+b+ab).
\]
Writing \(\bar x_i=x_i+I\), define the Golod group associated with these data by
\[
G=\langle 1+\bar x_1,\ldots,1+\bar x_r\rangle\leq 1+A.
\]
Here \(r\)-generated means generated by the displayed \(r\) elements; minimality of this generating set is not assumed.

\textbf{Statement (Kourovka Notebook Problem 9.76).} For every choice of \(p,k,r,I\) satisfying these assumptions and every \(g\in G\), the centralizer
\[
C_G(g)=\{h\in G:hg=gh\}
\]
is infinite.

The supplied record additionally reports constructions of Golod groups with infinite centre by A. V. Timofeenko and, independently by other methods in the 1990s, by L. Hammoudi. These are background claims about particular constructions, not additional hypotheses on \(G\). The record supplies no precise bibliographic reference for Hammoudi's work, and these historical claims have not been independently verified here.
\subsection{Short English statement}
For every Golod group formed from a finitely generated free associative algebra over a field of positive characteristic by taking a nil but non-nilpotent quotient as specified above, every element should commute with infinitely many elements of that group. Proposed by V. P. Shunkov in 1984; the supplied registry labels the problem open. This is a statement-only draft requiring human mathematical review.
\subsection{Sources}
\begin{itemize}
\item Local UnsolvedMath registry, supplied complete source record: permanent ID \(700609\), exact source problem code \texttt{catalogue entry}, title \emph{Kourovka Notebook Problem 9.76}. The supplied status is \texttt{open}; this label is not an independently verified literature assessment.
\item \emph{The Kourovka Notebook (Unsolved Problems in Group Theory)}, Problem 9.76, proposed by V. P. Shunkov, 9th issue (1984); also identified in the supplied record with the 21st issue edition (2026), \texttt{arXiv:1401.0300v48}. The source collection attributes its Issues 1--21 to editors E. I. Khukhro and V. D. Mazurov. The cited edition has not been independently consulted here.
\item E. S. Golod, \emph{Amer. Math. Soc. Transl. (2)} \textbf{48} (1965), 103--106, cited in the supplied statement for the algebraic construction. No article title is supplied.
\item A. V. Timofeenko, \emph{Math. Notes} \textbf{39}, no. 5 (1986), 353--355, cited in the supplied statement for constructions of Golod groups with infinite centre. No article title is supplied.
\item L. Hammoudi: the supplied statement attributes an independent construction by other methods to work in the 1990s, but supplies no title, publication venue, exact year, or bibliographic identifier.
\end{itemize}
\end{document}
